\documentclass[11pt,a4paper]{article}
\usepackage[margin=1in]{geometry}
\usepackage{booktabs}
\usepackage{longtable}
\usepackage{graphicx}
\usepackage{hyperref}
\usepackage{listings}
\usepackage{xcolor}
\usepackage{enumitem}
\usepackage{caption}
\usepackage{amsmath}
\usepackage{array}
\usepackage{tcolorbox}
\tcbuselibrary{breakable}

\newcolumntype{L}[1]{>{\raggedright\arraybackslash}p{#1}}

\hypersetup{
    colorlinks=true,
    linkcolor=blue,
    filecolor=magenta,
    urlcolor=cyan,
}

\tcbset{
    promptbox/.style={
        colback=blue!5!white,
        colframe=blue!75!black,
        title={#1},
        fonttitle=\bfseries,
        before skip=6pt,
        after skip=6pt,
        breakable,
    },
    notebox/.style={
        colback=orange!5!white,
        colframe=orange!70!black,
        title={#1},
        fonttitle=\bfseries,
        before skip=6pt,
        after skip=6pt,
        breakable,
    }
}

% Long unbreakable \texttt{} paths are the main source of overfull boxes;
% loosen line breaking instead of running into the margin.
\setlength{\emergencystretch}{3em}
\sloppy

\newcommand{\code}[1]{\texttt{\detokenize{#1}}}
\newcommand{\file}[1]{\textsl{#1}}
\newcommand{\sev}[1]{\textbf{#1}}

\title{Dead Code Report 02: Vulkan Engine \\
\large The Residual Set After the First Sweep --- Uncalled Methods, \\
Inline Accessors, Whole Units, and One Dormant Network Path}
\author{}
\date{\today}

\begin{document}

\maketitle

\begin{abstract}
This report re-audits the Vulkan engine at commit \texttt{9f0d76cb} (2026-10-24),
after the thirty-six \code{dead-code(...)} commits that followed report~01 and
after the three fixes derived from \code{perf_report_29}. The method is the
same linker-level relocation scan used in report~01 --- a symbol that is defined
but never a relocation target has no direct call, no address-taken use, and no
vtable entry --- followed by class-qualified textual verification, an
inline-accessor scan, and a cascade closure over everything only a dead caller
reaches.

The result is much smaller than report~01 and far more concentrated:
\textbf{74 non-trivial methods with zero call sites} (1{,}094 body lines),
\textbf{58 inline accessors and trivial methods} (179 lines), three whole units
(\code{EditableTexture} 212 lines, \code{Frustum} 67 lines,
\code{OctreeTopologyStats} 14 lines), and one dormant feature path --- the
\code{RemoteScene} mesh-fetch cluster (about 190 lines including members).
With declarations the deletable set is roughly \textbf{1{,}820 lines}.

Two findings are new relative to report~01. First, the perf work itself created
dead code: the batched-wakeup fix (L8) left \code{ThreadPool::enqueueDetached}
and \code{notifyOne} without callers. Second, the audit rejected an extension
that a fix agent had just added: the new per-phase fields on
\code{OctreeTopologyCompute::stats()} would have been dead on arrival because
\code{stats()} itself has no caller. The profiler output on stderr is the live
interface; the accessor and its \code{OctreeTopologyStats} header are deleted
instead.

Unlike report~01, no finding here touches a pipeline layout, a descriptor
binding, or a shader interface. Every deletion is CPU-side and
behaviour-preserving; the risk is link-level, not validation-level.
\end{abstract}

\tableofcontents
\newpage

\section{Scope and Methodology}
\label{sec:method}

\textbf{Snapshot.} HEAD \texttt{9f0d76cb} (\textit{debug: gate overlay
collection on the settings, instrument the walks}), 2026-10-24. The working
tree contains only the user's uncommitted \code{sdf/} work
(\code{TriangleDistanceFunction} and edits to the SVG/SDF primitives), which
this audit did not touch and which does not affect any finding below.
The tree is 47{,}418 C++ lines and 12{,}936 header lines across 521 tracked
\code{.cpp}/\code{.hpp}/\code{.tpp} files, excluding \code{third_party/},
\code{bin/}, and shaders. The debug build completes with zero warnings under
\code{-Wall -Wshadow}, so dead code remains invisible in the build log.

\textbf{Why a second report.} Report~01 (snapshot \texttt{6ba38f1}) found
129 methods plus roughly 5{,}000--6{,}500 lines across dead classes,
duplication, and feature paths. The 36 commits that followed removed the bulk
of that set. What remains is not a random sample of the old list: it is
(a) code that became dead \emph{after} report~01, including code killed by the
perf fixes; (b) methods the first scan could not separate because their names
are common or overloaded; (c) inline accessors, which have no symbol to scan;
and (d) one whole client feature path that report~01 classified as live.

\textbf{Techniques.}
\begin{enumerate}
    \item \textbf{Relocation scan.} For every project object in
    \code{bin/obj/debug}, the set of defined global text symbols
    (\code{nm -f posix}) was diffed against every relocation target
    (\code{objdump -r}) and undefined symbol across all objects. A symbol that
    is defined but never a relocation target has no direct call, no
    address-taken use, and no vtable/function-pointer entry anywhere in the
    project. The scan covered 1{,}009 unreferenced project method symbols
    before filtering; constructors, destructors, and compiler-generated
    symbols were discarded, leaving 356 call candidates and 136 unique
    class/method pairs.
    \item \textbf{Class-qualified textual verification.} Each candidate was
    re-checked with \code{grep -rn} over every \code{.cpp}/\code{.hpp} file,
    distinguishing declarations, definitions, log strings, and calls. This
    pass removes false positives from common names (e.g.\ \code{clear},
    \code{init}, \code{contains}) and confirms the relocation result.
    \item \textbf{Inline-accessor scan.} A parser over all headers collected
    every function with a body in a class definition and counted its
    occurrences in the whole tree. Fifty-eight definitions had exactly one
    occurrence --- the definition itself. Each was checked for pointer,
    macro, and string use before being listed.
    \item \textbf{Cascade closure.} For each dead method, its callees were
    inspected. A callee is deleted only if every caller is in the dead set
    (e.g.\ \code{Octree::fetch} $\rightarrow$ \code{getNodeAt(pos, level,
    simplification)}; \code{ImpostorCapture::captureAll} $\rightarrow$
    \code{capture}). Conversely, a method is \emph{not} deleted if a live
    caller exists, even when the relocation scan flagged it because the call
    is inlined or devirtualized.
    \item \textbf{Header-inclusion scan.} Every project header was searched as
    an \code{#include} target. Three headers are included but unused
    (\code{EditableTexture.hpp} by \code{VegetationRenderer.hpp}) or included
    by nobody (\code{Frustum.hpp}, \code{OctreeTopologyStats.hpp}).
\end{enumerate}

\textbf{What the scan cannot see.} Calls that exist only inside
\code{#ifdef USE_IMGUI} blocks in \emph{other} translation units are invisible
to the debug object tree only if the enclosing file is also excluded from the
debug build. Widget translation units are compiled in both modes (only the UI
runtime is gated), so the widget findings below come from real object files;
they were nevertheless re-checked textually in both preprocessor modes.
Virtual overrides and \code{std::function} targets are retained unless a
class-qualified search proves no call site exists. Line counts are
brace-matched body spans; they are estimates, not exact patch sizes.

\textbf{Not re-audited.} Shader binding cross-checks, the duplicate-block
detector, and the feature-path catalogue of report~01 are out of scope here.
The duplicate-block clusters that report~01 listed were either fixed or
explicitly deferred; re-running the detector would report the same deferred
items and is not evidence of new dead code.

\begin{notebox}{Self-audit of the perf-report fixes}
The three commits derived from \code{perf_report_29} were audited with the
same tools before this report was written.
\code{26f06f3e} (octree cache/fork/wakeup) made
\code{ThreadPool::enqueueDetached} and \code{notifyOne} dead, because
\code{Octree.cpp} now calls \code{enqueueDetachedNoNotify} and \code{notifyAll}
directly; both are deleted in \S\ref{sec:events}.
\code{7b32479a} (topology) added per-phase fields to
\code{OctreeTopologyCompute::stats()} and extended
\code{OctreeTopologyStats}. That extension was reverted before the commit:
\code{stats()} has no caller and no widget exposes it, so the new fields would
have been dead on arrival. The stderr profiler that \code{M7} actually asks for
is kept, and \code{stats()} plus its header are deleted here.
\code{9f0d76cb} (debug overlay gating) is live and unchanged.
\end{notebox}

\section{Summary of Findings}
\label{sec:summary}

\begin{center}\footnotesize
\begin{longtable}{p{0.50\textwidth}rr}
\toprule
Category & Count & Body lines \\
\midrule
\endfirsthead
\toprule
Category & Count & Body lines \\
\midrule
\endhead
\multicolumn{3}{l}{\textbf{A. Methods in live files with zero call sites}}\\
Services (\code{AudioAnalyzer}, \code{MaterialManager},
\code{OctreeTopologyCompute}, \code{TextureMixer},
\code{VideoTextureService}) & 15 & 143 \\
Core and resources (\code{VulkanApp}, \code{VulkanResourceManager},
\code{TextureArrayManager}) & 7 & 195 \\
Renderers (\code{SceneRenderer}, \code{SdfRenderer}, \code{VideoRenderer},
\code{WaterBackFaceRenderer}, \code{VegetationRenderer},
\code{ImpostorCapture}) & 17 & 209 \\
Events and space (\code{NunchukPublisher}, \code{EventManager},
\code{Octree}, \code{OctreeFile}) & 7 & 134 \\
Math (\code{BoundingBox}, \code{BoundingCube}, \code{Frustum}) & 9 & 68 \\
Utils and server (\code{GltfWriter}, \code{LocalScene},
\code{RemoteScene}, \code{ChunkProtocol}) & 12 & 133 \\
SDF (\code{SvgShape}, \code{SvgPolygonDistanceFunction}, \code{SdfScene}) & 4 & 81 \\
Application and widgets & 3 & 131 \\
\midrule
\textbf{Subtotal} & \textbf{74} & \textbf{1{,}094} \\
\midrule
\multicolumn{3}{l}{\textbf{B. Inline accessors and trivial methods}}\\
Header-defined getters, setters, and forwarders with one occurrence & 58 & 179 \\
\midrule
\multicolumn{3}{l}{\textbf{C. Whole units}}\\
\code{vulkan/resources/EditableTexture.{cpp,hpp}} (include-only) & 1 class & 212 \\
\code{math/geometry/Frustum.{cpp,hpp}} (never included) & 1 class & 21 \\
\code{services/OctreeTopologyStats.hpp} (only \code{stats()} uses it) & 1 struct & 14 \\
\midrule
\multicolumn{3}{l}{\textbf{D. Dormant feature path}}\\
\code{RemoteScene} client mesh fetch (\code{MSG_MESH_DATA}) & 1 path & $\sim$190 \\
\midrule
\multicolumn{3}{l}{\textbf{E. Deliberately kept}}\\
\code{IndirectRenderer::addMeshSlotted}/\code{uploadSlot} (latent slot API) & 2 & 296 \\
\code{SceneRenderer::bindSet0ForFrame} (descriptor-buffer fallback) & 1 & 73 \\
\midrule
\textbf{Deletable total (A+B+C+D)} & & $\approx$\textbf{1{,}820} \\
\bottomrule
\end{longtable}
\end{center}

The ``body lines'' column counts only the brace-matched method bodies, not the
declarations. Category~A's 74 declarations add roughly 110 lines. The
\code{Frustum} row in~C counts only the 21 lines outside its two methods, which
are already in~A (the unit is 67 lines in total); \code{EditableTexture}
includes its header. The deliberately kept items in~E are discussed in
\S\ref{sec:latent} and are not part of the total.

\section{Dead Functions}
\label{sec:deadfuncs}

Every method below has a defined symbol that is never a relocation target in
any project object \emph{and} no call-shaped textual use outside its
declaration and definition. The complete table with exact spans is
Appendix~\ref{app:full}.

\subsection{Services}
\label{sec:services}

\textbf{\code{AudioAnalyzer::setBandEdges(float, float)}} (4 lines,
\code{services/AudioAnalyzer.cpp:41--44}). Writes the bass/mid crossover
frequencies. The header comment says the defaults are ``kept configurable via
\code{setBandEdges()}'', but nothing configures them: \code{MusicWidget} calls
the window, spectrum-smoothing, amplitude-smoothing, and beat-sensitivity
setters, never this one. The comment is stale and should go with the method.

\textbf{\code{MaterialManager::destroy(VulkanApp*)}} (13 lines,
\code{vulkan/resources/MaterialManager.cpp:95--107}). Clears local handles and
the CPU cache but frees nothing: every buffer it names was registered with
\code{VulkanResourceManager} at creation, and the class destructor is
\code{= default}. It reads as a manual-ownership teardown that the codebase
does not have, and no caller exists.

\textbf{\code{OctreeTopologyCompute::destroy(VulkanApp*)}} (29 lines,
\code{services/OctreeTopologyCompute.cpp:137--165}). The buffers come from
\code{VulkanApp::createBuffer}, which registers each one with
\code{resources.addBufferVma} before returning; the pipeline, layout, and
descriptor pool are likewise registered. The destructor already prints the
profiler summary and clears the global instance. \code{destroy} is a second,
unused teardown path.

\textbf{\code{OctreeTopologyCompute::stats()}} (11 lines,
\code{services/OctreeTopologyCompute.cpp:783--793}). Returns the rolling
counters. It has no caller: no widget, no log path, and no test consumes it.
The live diagnostic is the \code{VULKAN_SETUP_DIAG} table printed from
\code{recordProfile}/\code{printProfileSummary}. Deleting \code{stats()} makes
\code{services/OctreeTopologyStats.hpp} dead (14 lines), so both go together.
See the self-audit box above: this is the accessor that a fix agent had just
extended.

\textbf{\code{TextureMixer::isLayerGenerationPending(uint32_t)}} (7 lines),
\textbf{\code{waitForLayerGeneration(VulkanApp*, uint32_t, uint64_t)}} (29
lines), and \textbf{\code{getGlobalInstance()}} (1 line,
\code{services/TextureMixer.cpp:32}). The wait loop is the only caller of the
singleton accessor, and nothing calls the wait loop. The static
\code{g_texture_mixer_instance} is written by the constructor and then never
read. Removing the pair also removes the only reason the singleton exists.

\textbf{\code{VideoTextureService}} (8 methods, 49 body lines). In the FFmpeg
branch: \code{backendCompiledIn} (7), \code{findAssetDirectory} (17),
\code{isSlotLoaded} (6), \code{videoArrayImage} (1), \code{videoArrayMips} (1),
\code{videoArrayHeight} (1), \code{slotOf} (5), \code{togglePlaying} (11).
The stub branch defines the same eight (\code{VideoTextureService.cpp:1955},
\code{:1957}, \code{:1962}, \code{:1964}, \code{:1973}, \code{:1977}).
Nothing queries whether FFmpeg is compiled in, nothing resolves an asset
directory through the service, and nothing reads the array image, mip count,
height, or slot mapping: the video renderer binds the service's texture array
through the service's own descriptor path. \code{togglePlaying} is not used by
\code{VideoTextureWidget}, which drives play/pause through the explicit
methods. Both branches must be edited together or the stub build breaks.

\subsection{Core and Resources}
\label{sec:core}

\textbf{\code{VulkanApp::generateMipmaps(VkImage, VkFormat, int32_t, int32_t,
uint32_t, uint32_t, uint32_t)}} (26 lines,
\code{vulkan/core/VulkanApp.cpp:1348--1373}). The live mip generator is
\code{ImpostorCapture::generateMipmaps}, which records into an existing
command buffer and is called twice from \code{recordCaptureViews}
(\code{ImpostorCapture.cpp:412--413}). The app-level helper predates it and has
no caller.

\textbf{\code{VulkanApp::recordTrackedLayoutForCommandBuffer}} (17 lines,
\code{vulkan/core/VulkanApp.cpp:3111--3127}). Layout tracking moved into the
recording helpers; this entry point is not called from any renderer.

\textbf{\code{VulkanApp::updateDescriptorSet(const std::vector<VkWriteDescriptorSet>&)}}
(37 lines, \code{vulkan/core/VulkanApp.cpp:3503--3539}). A full duplicate of
the \code{std::initializer_list} overload that follows it
(\code{:3541}). The list overload is live at
\code{widgets/RenderTargetsWidget.cpp:295} and
\code{vulkan/resources/SkySphere.cpp:54,151}; the vector overload is not. The
two bodies differ only in the container type and the same bad-handle guard, so
deleting the vector version removes 37 lines without touching a call site.

\textbf{\code{VulkanResourceManager::removeFramebuffer(VkFramebuffer)}} (1 line,
\code{vulkan/core/VulkanResourceManager.cpp:151}). Framebuffers are gone from
the dynamic-rendering path; the erase helper has no caller.

\textbf{\code{VulkanResourceManager::getSemaphoreMap()}} (1 line,
\code{:206}). Returns the internal map by const reference. No reader exists;
the resource-manager widget iterates its own snapshot.

\textbf{\code{TextureArrayManager::destroy(VulkanApp*)}} (21 lines,
\code{vulkan/resources/TextureArrayManager.cpp:178--198}) and
\textbf{\code{TextureArrayManager::load(...)}} (92 lines,
\code{:465--556}). \code{load} is the single-triple legacy path: it stages one
albedo/normal/bump/roughness/AO set into the next layer. The live bring-up path
is \code{loadTriples}, called twice from \code{MyApp.cpp:568} and
\code{:4506}. \code{destroy} has no caller because the arrays and views are
resource-manager owned; the live teardown is \code{releaseLayerViews} plus the
manager. This is the largest single-method deletion in the report.

\subsection{Renderers}
\label{sec:renderers}

\textbf{\code{SceneRenderer}} (3 methods, 27 lines).
\code{removeAllMeshesSolidOnly} (12 lines,
\code{vulkan/renderer/SceneRenderer.cpp:332--343}) and
\code{removeAllMeshesWaterOnly} (11 lines, \code{:345--355}) are reached only
by the inline wrappers \code{removeAllRegisteredMeshes} and
\code{removeAllTransparentMeshes} (\code{SceneRenderer.hpp:201--202}), which
are themselves uncalled; the live reset path is \code{resetSceneState} plus the
indirect renderers' own clears. \code{layerStateCount} (4 lines,
\code{:202--205}) is a debug accessor with no reader. The two
\code{GenPoolRef} accessors in the header (\code{:495--496}) exist only for
these methods and go with them.

\textbf{\code{SdfRenderer}} (7 methods, 42 lines). \code{setScene} (5,
\code{SdfRenderer.cpp:322--326}), \code{updateScene} (3, \code{:636--638}),
\code{setMarchParams} (6, \code{:679--684}), \code{setMarchRange} (7,
\code{:686--692}), \code{setMarchTiers} (9, \code{:714--722}),
\code{setSmokeWideContainer} (7, \code{:480--486}), and
\code{onSwapchainResized} (5, \code{:2142--2146}). The live scene path is
\code{pendingScene_}/\code{lavaScene_}/\code{smokeScene_} with
\code{applyScene}; the live march parameters are written through
\code{config()} by \code{RaymarchWidget}. \code{onSwapchainResized} is not an
override of \code{VulkanApp::onSwapchainResized} (different class) and is
never wired into the swapchain callback; \code{SceneRenderer} rebuilds SDF
targets through its own resize path. The inline \code{isRayMarchingEnabled}
and \code{hasScene} accessors (\code{SdfRenderer.hpp:146,278}) have no reader
and are part of the same group.

\textbf{\code{VideoRenderer::clear()}} (4 lines,
\code{vulkan/renderer/video/VideoRenderer.cpp:406--409}) and
\textbf{\code{colorLayout()}} (3 lines, \code{VideoRenderer.hpp:70--72}).
\code{clear} is a reset without a caller; teardown goes through
\code{cleanup}. \code{colorLayout} returns the tracked layout for a frame the
renderer already owns internally.

\textbf{\code{WaterBackFaceRenderer::init(VulkanApp*)}} (3 lines,
\code{vulkan/renderer/water/WaterBackFaceRenderer.cpp:14--16}). Not a virtual
override (the \code{Renderer} base declares only \code{cleanup},
\code{setCmdState}, and the destructor). Its whole body is
\code{appPtr = app}; the live \code{createRenderTargets} path sets the same
member (\code{:305}), so \code{init} is redundant rather than load-bearing.

\textbf{\code{VegetationRenderer::setBillboardArrayTextures(...)}} (18 lines,
\code{vulkan/renderer/vegetation/VegetationRenderer.cpp:1003--1020}). The live
descriptor refresh is the \code{TextureArrayManager} allocation listener
(\code{onTextureArraysReallocated} $\rightarrow$ \code{refreshVegDescriptors}),
which updates the existing set in place. Report~01 already identified the only
\code{MyApp} call to this setter as a no-op; that call was removed, and the
setter is now unreachable. The inline \code{setUseGpuPlacement},
\code{gpuPlacementEnabled}, and \code{isConsolidationDirty}
(\code{VegetationRenderer.hpp:115,116,163}) are vestigial accessors over live
state: \code{useGpuPlacement} and \code{vegConsolidationDirty} are read and
written by the render loop, so the accessors go and the members stay.

\textbf{\code{ImpostorCapture}} (4 methods, 115 lines).
\code{capture} (55 lines, \code{ImpostorCapture.cpp:454--508}),
\code{captureAll} (10, \code{:510--519}), \code{captureAsync} (38,
\code{:521--558}), and \code{updateTexDescSet} (12, \code{:1078--1089}).
These are the synchronous predecessors of the live asynchronous path:
\code{ImpostorService.cpp:47} calls \code{captureAllAsync}, which uses
\code{submitCaptureTypesAsync}, \code{completeCaptureType},
\code{recordCaptureViews}, and \code{writeCaptureCameraUbos}. \code{capture} is
called only by \code{captureAll}, and \code{captureAsync} by nobody.
\code{updateTexDescSet} is superseded by the descriptor writes inside the
async path. Deleting the four leaves the entire live capture path intact.

\subsection{Events and Space}
\label{sec:events}

\textbf{\code{NunchukPublisher::aButtonPressed()}} (7 lines,
\code{events/NunchukPublisher.cpp:792--798}) and the \code{aPrev} member
(\code{NunchukPublisher.hpp:74}). \code{GamepadPublisher} has an identically
named, live pair (\code{GamepadPublisher.cpp:59--63}, called from
\code{RadialMenuHandler.cpp:192}); the Nunchuk one is never polled. The
duplicate name is why a name-only scan cannot separate them.

\textbf{\code{EventManager::unsubscribe(HandlerPtr)}} (5 lines,
\code{events/EventManager.cpp:21--25}). Handler lists are cleared wholesale on
teardown; nothing unsubscribes a single handler.

\textbf{\code{Octree::fetch(...)}} (10 lines, \code{space/Octree.cpp:218--227}),
\textbf{\code{Octree::getNodeAt(const glm::vec3&, int, bool)}} (23 lines,
\code{:145--167}), and \textbf{\code{Octree::getNodeAt(const glm::vec3&,
bool)}} (21 lines, \code{:169--189}). \code{fetch} calls the three-argument
overload; the two-argument overload has no caller at all. The live spatial
query path is \code{intersect} and the chunk-manager lookups. This is a clean
three-method cascade: delete the entry point and both query overloads together.

\textbf{\code{OctreeFile::load(std::string, float)}} (33 lines,
\code{space/OctreeFile.cpp:61--93}) and \textbf{\code{save(std::string,
float)}} (35 lines, \code{:95--129}). The live serialization path is
\code{readFromStream}/\code{writeToStream}, used by \code{LocalScene.cpp:623,
669--670,697} and \code{MyApp.cpp:4884,4972}. \code{load}/\code{save} add file
naming and zlib framing on top and are called by nobody.

\subsection{Math}
\label{sec:math}

\textbf{\code{BoundingBox::setMaxX/setMaxZ/setMinX/setMinZ}} (4 $\times$ 3
lines, \code{math/geometry/BoundingBox.cpp:48--70}). The whole-vector
\code{setMax}/\code{setMin} are live. The Y-axis pair is also live:
\code{utils/MainSceneLoader.hpp:456--457} clamps the water box with
\code{setMaxY(0)} and \code{setMinY(...)}. Only the X/Z per-axis setters are
unused. This is exactly the kind of asymmetry that a name-based sweep would get
wrong in both directions.

\textbf{\code{BoundingCube::getChildCenter(int)}} (4 lines,
\code{math/geometry/BoundingCube.cpp:59--62}),
\textbf{\code{BoundingCube::intersects(const Ray&)}} (3 lines, \code{:65--67}),
and \textbf{\code{BoundingCube::setMax(glm::vec3)}} (3 lines, \code{:69--71}).
The box/box \code{intersects} inherited from \code{AbstractBoundingBox} is live;
the ray overload is a separate non-override with no caller. \code{getChild}
already returns the child cube, so \code{getChildCenter} has no user.
\code{setMax} is the cube-specific setter; the live path uses
\code{setLength} and the base-class extents.

\textbf{\code{Ray::getOrigin()}} (1 line, \code{math/geometry/Ray.hpp:12}).
\code{Ray} itself is live (the Nunchuk aim raycast at
\code{NunchukPublisher.cpp:700}); only the accessor is unused.

\textbf{\code{Frustum}} (whole unit, 67 lines). No file includes
\code{Frustum.hpp}; the constructor and \code{test} are defined in
\code{Frustum.cpp} and never referenced. \code{ContainmentType} is used by
other geometry types, but the frustum class is not part of any visibility
path. See \S\ref{sec:units}.

\subsection{Utils and Server}
\label{sec:utils}

\textbf{\code{GltfWriter::vertexCount()}} and \textbf{\code{triangleCount()}}
(6 lines each, \code{utils/GltfWriter.cpp:93--105}). Statistics accessors for
the exporter; the exporter logs node and mesh counts through other paths.

\textbf{\code{LocalScene::getOpaqueOctree()}} and
\textbf{\code{getTransparentOctree()}} (both overloads, 14 lines total,
\code{utils/LocalScene.cpp:297--310}). The live accessor is
\code{getLayerOctree(Layer)}; \code{MyApp} and the renderers use the layer
API. The four typed getters are convenience wrappers with no user.

\textbf{\code{RemoteScene::fetchMesh}}, \textbf{\code{scheduleRetry}},
\textbf{\code{buildGeometry}}, and \textbf{\code{killStub}} (96 lines total)
are the visible part of the dormant mesh-fetch path; they are described
together with the receiver branch and members in \S\ref{sec:remote}.

\textbf{\code{chunkproto::buildRequestMesh(uint64_t)}} (8 lines,
\code{server/ChunkProtocol.cpp:90--97}) and
\textbf{\code{buildRequestAllMeshes()}} (3 lines, \code{:99--101}). These are
the client-side request builders. They become unreachable when the fetch path
goes; the server-side handlers (\code{WebServer.cpp:750,758}) and the message
constants stay, because other clients speak the same protocol.

\textbf{\code{WsClient::sendBinary(const uint8_t*, size_t)}} (3 lines,
\code{utils/WsClient.cpp:198--200}). A forwarding overload over the
\code{std::vector} version. Its only caller was \code{fetchMesh}; the live
snapshot request and every other send use the vector overload.

\textbf{\code{WebServer::textureRoot()}} (1 line, \code{server/WebServer.hpp:50}).
Accessor for the texture root; no reader.

\subsection{SDF}
\label{sec:sdf}

\textbf{\code{SvgShape::rgbToHsvDegrees(const glm::vec3&)}} (14 lines,
\code{sdf/SvgShape.cpp:492--505}) and \textbf{\code{SvgShape::contains(const
glm::vec2&)}} (3 lines, \code{:507--509}). The HSV conversion was for brush
tinting on the SVG path; the live tint path is the Nunchuk colour subpage,
which has its own conversion. \code{contains} is a point-in-polygon test; the
distance function uses \code{localPoint} and its own winding logic.

\textbf{\code{SvgPolygonDistanceFunction::localProfilePoint(const glm::vec3&)}}
(6 lines, \code{sdf/SvgPolygonDistanceFunction.cpp:88--93}). Normalized
SVG-space projection; \code{localPoint} is the live one.

\textbf{\code{SdfScene::createFireDemo(uint32_t, glm::vec3, float)}} (58
lines, \code{sdf/types/SdfScene.cpp:363--420}). A demo fire ring built from
evenly spaced anchors. The live fire path is
\code{createFireFromAnchors}, which \code{SdfRenderer} calls with anchors from
the scene (\code{SdfRenderer.cpp:429,862}). The demo has no loader and no
caller; its own comment notes that it must be kept visually in sync with
\code{createFireFromAnchors}, which is maintenance cost for dead code.

\subsection{Application and Widgets}
\label{sec:app}

\textbf{\code{MyApp::dispatchSolidEvents()}} (3 lines,
\code{MyApp.cpp:5195--5197}) and \textbf{\code{dispatchLiquidEvents()}} (8
lines, \code{:5199--5206}). Both are thin wrappers over
\code{dispatchLayerEvents(LAYER_OPAQUE/LAYER_TRANSPARENT)}. The live pump is
\code{dispatchAllEvents}, which snapshots the collector count and dispatches
every layer; nothing calls the per-layer pair.

\textbf{\code{RenderTargetsWidget::cleanup()}} (120 lines,
\code{widgets/RenderTargetsWidget.cpp:662--781}). A near-copy of the
destructor's owned-descriptor sweep, extended with the linear-target image
teardown and the linearization pipeline, layout, descriptor set, and sampler
destruction. It is never called. The live paths already cover the work:
\code{destroyLinearTargets()} (called from \code{setFrameInfo} on resize)
destroys the images, views, and allocations, and the pipeline, layout, and
sampler are registered with \code{VulkanResourceManager}. The destructor
retains the descriptor removal. Deleting \code{cleanup} also removes a
duplicated \code{removeOwnedDesc(waterDepthLinearDescriptor, ...)} line: the
second call is harmless because the first nulls the handle, but it is a copy
error that only exists in dead code.

\textbf{\code{MusicWidget::isPlaying()}} (1 line,
\code{widgets/MusicWidget.hpp:44}), \textbf{\code{OctreeExplorerWidget::
setExportBlenderCallback(...)}} (3 lines,
\code{widgets/OctreeExplorerWidget.hpp:21--23}), and
\textbf{\code{ImGuiHelpers::openPtr(...)}} (1 line,
\code{widgets/components/ImGuiHelpers.hpp:22}). The music widget drives the
analyzer directly; the octree explorer never registers an export callback (the
Blender export is started from \code{MyApp}); the pointer-opening helper was
replaced by the reference overload that is actually used.

\section{Whole Units}
\label{sec:units}

\textbf{\code{vulkan/resources/EditableTexture.{cpp,hpp}} (212 lines).} The
class is the old per-layer editable CPU texture: allocate, upload, draw an
ImGui editor, invalidate the descriptor on pool rebuild. It was replaced by the
\code{TextureArrayManager} arrays, which own the albedo/normal/bump/roughness/
AO layers and expose them through \code{getImTexture}. The only remaining
reference is \code{#include "../../resources/EditableTexture.hpp"} at
\code{vulkan/renderer/vegetation/VegetationRenderer.hpp:6}; the class is never
named in that header or its source. Deleting the include also removes an
unnecessary \code{<imgui.h>} dependency from the vegetation renderer header.
The files are compiled (they are in \code{vulkan/resources/*.cpp}), so the
symbols exist and the relocation scan sees them; they are simply unreachable.

\textbf{\code{math/geometry/Frustum.{cpp,hpp}} (67 lines).} The header is
included by nobody. The class has a matrix constructor that extracts six planes
and a \code{test} that classifies a bounding box. No renderer, culler, or
widget uses it; visibility is decided by the octree and the indirect renderer's
own bounds. Report~01 listed \code{Frustum} as an orphan candidate but did not
delete it; this report closes it.

\textbf{\code{services/OctreeTopologyStats.hpp} (14 lines).} A plain struct of
rolling counters, included only by \code{OctreeTopologyCompute.hpp} and
returned only by the uncalled \code{stats()}. The live profiler prints its own
per-window table and does not use the struct. Delete the header, the
\code{stats()} declaration, and the include together.

\section{Inline Accessors and Vestigial State}
\label{sec:inline}

Inline functions have no out-of-line symbol, so the relocation scan cannot see
them. The header parser found 58 definitions whose name occurs exactly once in
the tree --- the definition itself. They fall into four groups.

\begin{center}\footnotesize
\begin{longtable}{p{0.42\textwidth}p{0.40\textwidth}r}
\toprule
File & Symbol & Lines \\
\midrule
\endfirsthead
\toprule
File & Symbol & Lines \\
\midrule
\endhead
\code{services/AudioAnalyzer.hpp:47} & \code{windowType} & 1 \\
\code{services/AudioAnalyzer.hpp:52} & \code{spectrumSmoothing} & 1 \\
\code{services/AudioAnalyzer.hpp:53} & \code{amplitudeSmoothing} & 1 \\
\code{services/AudioAnalyzer.hpp:58} & \code{beatSensitivity} & 1 \\
\code{services/AudioAnalyzer.hpp:77} & \code{spectrumData} & 1 \\
\code{services/AudioAnalyzer.hpp:78} & \code{spectrumBins} & 1 \\
\code{services/AudioAnalyzer.hpp:90} & \code{waveformCapacity} & 1 \\
\code{services/ParticlePhysicsSim.hpp:107} & \code{regionNodeCount} & 1 \\
\code{sdf/SvgPolygonDistanceFunction.hpp:35} & \code{halfExtents} & 1 \\
\code{utils/MainSceneLoader.hpp:83} & \code{shapeIndex} & 3 \\
\code{utils/LayerRendererType.hpp:15} & \code{layerRendererName} & 7 \\
\code{math/TextureType.hpp:25} & \code{textureTypeName} & 10 \\
\code{math/geometry/Ray.hpp:12} & \code{getOrigin} & 1 \\
\code{widgets/MusicWidget.hpp:44} & \code{isPlaying} & 1 \\
\code{widgets/OctreeExplorerWidget.hpp:21} & \code{setExportBlenderCallback} & 3 \\
\code{widgets/components/ImGuiHelpers.hpp:22} & \code{openPtr} & 1 \\
\code{vulkan/core/VulkanApp.hpp:213} & \code{getCurrentFrameFence} & 4 \\
\code{vulkan/renderer/SceneRenderer.hpp:201} & \code{removeAllRegisteredMeshes} & 1 \\
\code{vulkan/renderer/SceneRenderer.hpp:202} & \code{removeAllTransparentMeshes} & 1 \\
\code{vulkan/renderer/SceneRenderer.hpp:495} & \code{solidGenPoolRef} & 1 \\
\code{vulkan/renderer/SceneRenderer.hpp:496} & \code{waterGenPoolRef} & 1 \\
\code{vulkan/renderer/DescriptorUpdateStats.hpp:19} & \code{updateCalls} & 3 \\
\code{vulkan/renderer/SlotAllocator.hpp:115} & \code{updateCounts} & 11 \\
\code{vulkan/renderer/SlotAllocator.hpp:128} & \code{getSlot} & 7 \\
\code{vulkan/renderer/SlotAllocator.hpp:149} & \code{totalVertexCapacity} & 3 \\
\code{vulkan/renderer/SlotAllocator.hpp:154} & \code{totalIndexCapacity} & 3 \\
\code{vulkan/renderer/SceneQueues.hpp:28} & \code{getGeometryQueue} & 1 \\
\code{vulkan/renderer/CommandBufferState.hpp:34} & \code{bindComputePipeline} & 5 \\
\code{vulkan/renderer/CommandBufferState.hpp:58} & \code{bindGraphicsPipelineUnchecked} & 4 \\
\code{vulkan/renderer/DescriptorBuffer.hpp:36} & \code{bindingOffset} & 9 \\
\code{vulkan/renderer/DescriptorBufferHelper.hpp:22} & \code{bindingInfo} & 8 \\
\code{vulkan/renderer/solid/SolidRenderer.hpp:92} & \code{activeGraphicsPipeline} & 6 \\
\code{vulkan/renderer/rt/RayTracingResources.hpp:158} & \code{getRTSetLayout} & 1 \\
\code{vulkan/renderer/sdf/SdfRenderer.hpp:146} & \code{isRayMarchingEnabled} & 1 \\
\code{vulkan/renderer/sdf/SdfRenderer.hpp:278} & \code{hasScene} & 1 \\
\code{vulkan/renderer/vegetation/VegetationRenderer.hpp:115} & \code{setUseGpuPlacement} & 1 \\
\code{vulkan/renderer/vegetation/VegetationRenderer.hpp:116} & \code{gpuPlacementEnabled} & 1 \\
\code{vulkan/renderer/vegetation/VegetationRenderer.hpp:163} & \code{isConsolidationDirty} & 1 \\
\code{vulkan/renderer/video/VideoRenderer.hpp:70} & \code{colorLayout} & 3 \\
\code{vulkan/renderer/water/WaterRenderer.hpp:153} & \code{rtShadingEnabled} & 1 \\
\code{vulkan/renderer/water/WaterRenderer.hpp:189} & \code{rtProfilingEnabled} & 1 \\
\code{vulkan/renderer/indirect/IndirectRenderer.hpp:64} & \code{hasUploadManager} & 1 \\
\code{vulkan/renderer/indirect/IndirectRenderer.hpp:282} & \code{getVegTableBuffer} & 1 \\
\code{vulkan/renderer/indirect/IndirectRenderer.hpp:283} & \code{isVegetationCullEnabled} & 1 \\
\code{vulkan/renderer/indirect/IndirectRenderer.hpp:286} & \code{getVegDummyBuffer} & 1 \\
\code{vulkan/renderer/indirect/IndirectRenderer.hpp:304} & \code{getIndirectBuffer} & 1 \\
\code{vulkan/renderer/indirect/IndirectRenderer.hpp:305} & \code{getBoundsBuffer} & 1 \\
\code{vulkan/renderer/indirect/IndirectRenderer.hpp:306} & \code{getCompactBuffer} & 1 \\
\code{vulkan/renderer/indirect/IndirectRenderer.hpp:314} & \code{getLastVisibleCount} & 1 \\
\code{vulkan/renderer/indirect/IndirectRenderer.hpp:315} & \code{getComputeDescriptorSetLayout} & 1 \\
\code{vulkan/renderer/indirect/IndirectRenderer.hpp:342} & \code{visitActiveMeshInfos} & 7 \\
\code{vulkan/renderer/indirect/IndirectRenderer.hpp:360} & \code{copyRTGeometrySpans} & 32 \\
\code{vulkan/renderer/indirect/IndirectRenderer.hpp:479} & \code{getVertexBufferSize} & 1 \\
\code{vulkan/renderer/indirect/IndirectRenderer.hpp:480} & \code{getIndexBufferSize} & 1 \\
\code{vulkan/includes/shader/vertex_layouts.hpp:14} & \code{defaultBindings} & 3 \\
\code{vulkan/includes/shader/vertex_layouts.hpp:18} & \code{vertexAndInstanceBindings} & 6 \\
\code{vulkan/includes/shader/vertex_layouts.hpp:37} & \code{defaultAttributesWithInstance} & 5 \\
\code{server/WebServer.hpp:50} & \code{textureRoot} & 1 \\
\midrule
\textbf{Total} & \textbf{58 symbols} & \textbf{179} \\
\bottomrule
\end{longtable}
\end{center}

Three caveats apply to this table.

\begin{enumerate}
    \item \textbf{Grouped families.} The \code{IndirectRenderer} getters
    (\code{getIndirectBuffer}, \code{getBoundsBuffer}, \code{getCompactBuffer},
    \code{getLastVisibleCount}, \code{getComputeDescriptorSetLayout},
    \code{hasUploadManager}, \code{getVegTableBuffer}, \code{getVegDummyBuffer},
    \code{isVegetationCullEnabled}, \code{visitActiveMeshInfos},
    \code{copyRTGeometrySpans}) are a deliberate inspection surface: they exist
    so the resource-manager widget and external tools can read the renderer's
    buffers. They have no in-tree caller, but deleting them removes the only
    supported way to inspect those handles without exposing members. The
    conservative choice is to keep this family and delete only the four
    \code{SlotAllocator} helpers and the two size getters, which duplicate
    public capacity fields. The prompt in \S\ref{sec:prompt} takes that
    conservative branch.
    \item \textbf{State, not just code.} \code{DescriptorUpdateStats::updateCalls}
    is written by the update path and never read; deleting its accessor is only
    half the change, and the counter member must go with it. The three
    \code{VegetationRenderer} accessors are not in that class: the backing
    members \code{useGpuPlacement} and \code{vegConsolidationDirty} are live
    (the render loop reads and writes both), so only the accessors are deleted
    and the members stay.
    \item \textbf{Name collisions.} \code{getSlot} is not
    \code{ChunkManager::getSlotIndex}; \code{textureTypeName} is not
    \code{layerRendererName}. Each was checked with its class qualifier before
    being listed.
\end{enumerate}

\section{Dead Feature Path: RemoteScene Mesh Fetch}
\label{sec:remote}

The remote scene speaks the chunk protocol in \code{server/ChunkProtocol.hpp}.
The server can answer \code{MSG_REQUEST_MESH} and
\code{MSG_REQUEST_ALL_MESHES} with \code{MSG_MESH_DATA}
(\code{server/WebServer.cpp:750--764}, \code{ChunkProtocol.cpp:190}). The
client never asks. The only sender of those requests was
\code{RemoteScene::fetchMesh}, which has no caller anywhere in the tree, so the
receiver's \code{MSG_MESH_DATA} branch can never see a reply to a request it
made. The snapshot path (\code{MSG_REQUEST_ALL} $\rightarrow$
\code{START/META/BATCH/END}) and the live batch path are separate and remain
live.

The dormant cluster is larger than the four methods because it owns state:

\begin{center}\footnotesize
\begin{tabular}{p{0.52\textwidth}rr}
\toprule
Part & Lines & Note \\
\midrule
\code{fetchMesh} & 23 & request, 30\,s wait, reply move \\
\code{scheduleRetry} & 25 & 5\,s throttle via \code{lastRetryMs_} \\
\code{buildGeometry} & 39 & normals, triplanar UVs, brush ids \\
\code{killStub} & 9 & delete path flips \code{alive} directly \\
\code{MSG_MESH_DATA} branch & 33 & receiver-side reply matching \\
\code{triplanarPlane}/\code{triplanarMapping} & 18 & used only by \code{buildGeometry} \\
\code{CacheKey}/\code{CacheKeyHash} + \code{meshCache_} & 14 & cache is erased, never inserted \\
\code{lastRetryMs_} + \code{cacheMutex_} & 2 & only the two dead paths touch them \\
\code{fetchMutex_}/\code{fetchCv_}/\code{fetchPending_}/ & 8 & reply slot and wakeups \\
\code{fetchId_}/\code{fetchReply_}/\code{fetchReady_} & & \\
\code{disconnectLocked}/\code{receiverLoop} wakeups & 6 & only wake the fetch wait \\
\midrule
\textbf{Total} & $\approx$\textbf{190} & \\
\bottomrule
\end{tabular}
\end{center}

Two details make the path look alive to a superficial check. First,
\code{meshCache_} appears in the delete handler
(\code{RemoteScene.cpp:483--485}), where it prunes entries for a deleted
chunk --- but nothing ever inserts an entry, so the loop iterates an always-empty
map. Second, \code{disconnectLocked} and \code{receiverLoop} both call
\code{fetchCv_.notify_all()}; those wakeups exist only to unblock
\code{fetchMesh}'s 30-second wait. The delete path already flips
\code{stub->alive = false} itself (\code{:422,480}), so \code{killStub} is a
third copy of the same operation.

\textbf{What is deleted.} \code{fetchMesh}, \code{scheduleRetry},
\code{buildGeometry}, \code{killStub}, the \code{MSG_MESH_DATA} branch, the
two triplanar helpers, \code{CacheKey}/\code{CacheKeyHash}, \code{meshCache_},
\code{cacheMutex_}, \code{lastRetryMs_}, \code{fetchMutex_}, \code{fetchCv_},
\code{fetchPending_}, \code{fetchId_}, \code{fetchReply_},
\code{fetchReady_}, the fetch lines in \code{disconnectLocked} and
\code{receiverLoop}, the cache-prune loop in the delete handler, the
\code{<condition_variable>} include if nothing else needs it, and the two
\code{chunkproto} request builders that only \code{fetchMesh} used
(\code{buildRequestMesh}, \code{buildRequestAllMeshes}).

\textbf{What is kept.} All message constants, including
\code{MSG_REQUEST_MESH}, \code{MSG_REQUEST_ALL_MESHES}, and
\code{MSG_MESH_DATA}; the server handlers; the server's
\code{buildMeshData}; \code{sendMutex_} (the live snapshot request at
\code{RemoteScene.cpp:331} uses it); \code{findByNode}; \code{reemitLayer}; and
the delete path. This is a client-feature deletion, not a protocol deletion:
the server keeps serving mesh requests to any other client, and the protocol
header remains the single source of truth for the wire format.

\section{Deliberately Kept: Latent API}
\label{sec:latent}

Two uncalled APIs are \emph{not} in the delete set, and the distinction matters
for the prompt.

\textbf{\code{IndirectRenderer::addMeshSlotted}} (96 lines,
\code{IndirectRenderer.cpp:2241--2336}) and
\textbf{\code{uploadSlot}} (200 lines, \code{:2736--2935}). A complete
slotted-upload interface: allocate a slot, upload vertex/index spans, publish
the slot. The streaming engine currently goes through the chunk-manager path
and the upload manager; the slotted path is the target of the in-progress
migration described in the renderer's own comments. It is reachable by design
and would be re-added immediately if deleted. Keep, and mark with a comment
naming the migration, so the next audit does not re-litigate it.

\textbf{\code{SceneRenderer::bindSet0ForFrame}} (73 lines,
\code{SceneRenderer.cpp:1384--1456}). The descriptor-buffer fallback for set~0
when the main layout carries \code{DESCRIPTOR_BUFFER_BIT}. Report~01 already
classified it as intentionally uncalled; the layout still has the conditional
comment at \code{SceneRenderer.hpp:466}, so the API is part of a supported
mode, not an accident. Keep.

The difference from the delete set is not size but intent: a latent API has a
documented migration or a supported mode behind it, and deleting it would
remove a capability. A dead method is a superseded duplicate or a leftover
entry point, and deleting it removes only maintenance cost. The
\code{RenderTargetsWidget::cleanup} case shows the test: it duplicates the
destructor and \code{destroyLinearTargets} with no capability that the live
paths lack, so it is dead, not latent.

\section{Cascades and Deletion Order}
\label{sec:cascade}

The deletions are grouped so that each commit leaves the tree link-clean and
each group's justification is local.

\begin{enumerate}
    \item \textbf{Services.} \code{AudioAnalyzer::setBandEdges};
    \code{MaterialManager::destroy};
    \code{OctreeTopologyCompute::{destroy,stats}} plus
    \code{OctreeTopologyStats.hpp};
    \code{TextureMixer::{isLayerGenerationPending,waitForLayerGeneration,
    getGlobalInstance}} plus the static pointer;
    \code{VideoTextureService}'s eight methods in both branches.
    \item \textbf{Core and resources.} \code{VulkanApp::{generateMipmaps,
    recordTrackedLayoutForCommandBuffer, updateDescriptorSet(vector)}};
    \code{VulkanResourceManager::{removeFramebuffer,getSemaphoreMap}};
    \code{TextureArrayManager::{destroy,load}}.
    \item \textbf{Renderers.} \code{SceneRenderer}'s two remove methods and
    \code{layerStateCount} plus the wrappers and pool refs;
    \code{SdfRenderer}'s seven methods plus the two inline accessors;
    \code{VideoRenderer::{clear,colorLayout}};
    \code{WaterBackFaceRenderer::init};
    \code{VegetationRenderer::setBillboardArrayTextures} plus the three inline
    flags; \code{ImpostorCapture::{capture,captureAll,captureAsync,
    updateTexDescSet}}.
    \item \textbf{Events, space, math.} \code{NunchukPublisher::
    aButtonPressed} plus \code{aPrev}; \code{EventManager::unsubscribe};
    \code{Octree::fetch} and both \code{getNodeAt} overloads;
    \code{OctreeFile::{load,save}};
    \code{BoundingBox}'s four X/Z setters;
    \code{BoundingCube::{getChildCenter,intersects(Ray),setMax}};
    \code{Ray::getOrigin}; the \code{Frustum} unit.
    \item \textbf{Utils, server, SDF.} \code{GltfWriter}'s two counters;
    \code{LocalScene}'s four typed getters; the \code{RemoteScene} fetch
    cluster; \code{chunkproto}'s two request builders;
    \code{WsClient::sendBinary(const uint8_t*, size_t)};
    \code{WebServer::textureRoot}; \code{SvgShape}'s two methods;
    \code{SvgPolygonDistanceFunction::localProfilePoint};
    \code{SdfScene::createFireDemo}.
    \item \textbf{Application and widgets.} \code{MyApp}'s two dispatch
    wrappers; \code{RenderTargetsWidget::cleanup}; the three widget accessors;
    the conservative subset of the inline table (the four \code{SlotAllocator}
    helpers and the two \code{IndirectRenderer} size getters).
    \item \textbf{Perf-fix residue.} \code{ThreadPool::enqueueDetached} and
    \code{notifyOne}, which the L8 fix made dead.
\end{enumerate}

Ordering inside a group matters in four places. \code{Octree::fetch} must go
with both \code{getNodeAt} overloads: deleting \code{fetch} first leaves the
three-argument overload with no caller and the two-argument overload dead from
the start. \code{ImpostorCapture::capture} must go with \code{captureAll}
(its only caller) but \emph{not} with \code{captureAllAsync}, which is live.
\code{OctreeTopologyCompute::stats} must go with
\code{OctreeTopologyStats.hpp}, or the header becomes an orphan. The
\code{RemoteScene} methods must go with the receiver branch and the members,
or the branch keeps reading a reply slot nobody waits on.

\section{Verification Plan}
\label{sec:verify}

\begin{enumerate}
    \item \code{make -j8 debug} after each group. The debug build compiles
    every project translation unit, including \code{widgets/*.cpp},
    \code{server/*.cpp}, and \code{space/*.cpp}; a deletion that breaks a
    caller fails here. Output is teed to \code{logs/build.log} by the make
    target.
    \item \code{grep -rn} for each removed name over \code{.cpp}/\code{.hpp}/
    \code{.tpp}; the only allowed hits are comments that mention the old API,
    and those comments should be updated or removed in the same commit.
    \item Link check: \code{make -j8 debug} links \code{bin/app}. The server
    binary is not built by the agent (release target), but the deleted
    server-side symbols are compiled into the app's objects and their absence
    of callers is textual, so the server link risk is nil.
    \item No runtime validation is performed as part of this audit. Every
    deletion is CPU-side: no pipeline layout, descriptor binding, shader
    interface, or synchronization object changes. The usual
    \code{VULKAN_GPU_ASSISTED=1} run is therefore a formality for this set,
    not a correctness gate. If a later commit mixes these deletions with
    renderer-state changes, that run becomes mandatory again.
    \item Widget paths: the debug build compiles the widget translation units
    but does not run the ImGui UI. The four widget deletions were verified
    textually in both preprocessor modes; a release run remains the user's
    check, not the agent's.
\end{enumerate}

\section{Proposed Prompt}
\label{sec:prompt}

The prompt below is written to be pasted into a coding agent in this
repository. It is directory-staged, forbids mixing the deletions with
behaviour changes, and names the exact symbols so the agent does not re-run the
audit from scratch.

\begin{tcolorbox}[promptbox={Prompt --- Dead-code sweep 02}]
Goal: delete the residual dead code identified in
\code{docs/dead_code_02.tex} at \code{/home/hdlrocha/vulkan-engine}. Read
\code{AGENTS.md} first. Never modify \code{third_party/}. Do not mix these
deletions with behaviour changes, refactors, or formatting. Do not touch the
uncommitted \code{sdf/} work.

Baseline: HEAD \code{9f0d76cb}. For every symbol, confirm with a
class-qualified \code{grep -rn} before deleting, and remove the header
declaration with the definition. If a symbol turns out to have a live caller,
leave it and report why.

Group 1 --- services. Delete \code{AudioAnalyzer::setBandEdges};
\code{MaterialManager::destroy}; \code{OctreeTopologyCompute::destroy} and
\code{stats()} plus \code{services/OctreeTopologyStats.hpp} and its include;
\code{TextureMixer::isLayerGenerationPending}, \code{waitForLayerGeneration},
\code{getGlobalInstance}, and \code{g_texture_mixer_instance};
\code{VideoTextureService::backendCompiledIn}, \code{findAssetDirectory},
\code{isSlotLoaded}, \code{videoArrayImage}, \code{videoArrayMips},
\code{videoArrayHeight}, \code{slotOf}, \code{togglePlaying} in \emph{both}
the FFmpeg and stub branches.

Group 2 --- core and resources. Delete \code{VulkanApp::generateMipmaps},
\code{recordTrackedLayoutForCommandBuffer}, and the
\code{std::vector<VkWriteDescriptorSet>} overload of
\code{updateDescriptorSet} (keep the \code{std::initializer_list} overload);
\code{VulkanResourceManager::removeFramebuffer} and \code{getSemaphoreMap};
\code{TextureArrayManager::destroy} and \code{load} (keep
\code{loadTriples}).

Group 3 --- renderers. Delete \code{SceneRenderer::removeAllMeshesSolidOnly},
\code{removeAllMeshesWaterOnly}, \code{layerStateCount}, the wrappers
\code{removeAllRegisteredMeshes}/\code{removeAllTransparentMeshes}, and the
two \code{GenPoolRef} accessors; \code{SdfRenderer::setScene},
\code{updateScene}, \code{setMarchParams}, \code{setMarchRange},
\code{setMarchTiers}, \code{setSmokeWideContainer},
\code{onSwapchainResized}, plus the inline \code{isRayMarchingEnabled} and
\code{hasScene}; \code{VideoRenderer::clear} and \code{colorLayout};
\code{WaterBackFaceRenderer::init};
\code{VegetationRenderer::setBillboardArrayTextures} plus the inline
\code{setUseGpuPlacement}, \code{gpuPlacementEnabled}, and
\code{isConsolidationDirty} accessors (the backing members are live; keep
them);
\code{ImpostorCapture::capture}, \code{captureAll}, \code{captureAsync}, and
\code{updateTexDescSet}. Keep \code{captureAllAsync} and its helpers.

Group 4 --- events, space, math. Delete \code{NunchukPublisher::aButtonPressed}
and \code{aPrev} (do not touch \code{GamepadPublisher}'s live pair);
\code{EventManager::unsubscribe}; \code{Octree::fetch} and both
\code{getNodeAt} overloads; \code{OctreeFile::load} and \code{save} (keep
\code{readFromStream}/\code{writeToStream}); \code{BoundingBox}'s
\code{setMaxX}, \code{setMaxZ}, \code{setMinX}, \code{setMinZ} (keep the Y
setters, which the water box uses); \code{BoundingCube::getChildCenter},
\code{intersects(const Ray&)}, and \code{setMax}; \code{Ray::getOrigin};
delete \code{math/geometry/Frustum.{cpp,hpp}}.

Group 5 --- utils, server, SDF. Delete \code{GltfWriter::vertexCount} and
\code{triangleCount}; \code{LocalScene}'s four typed octree getters; the
\code{RemoteScene} fetch path exactly as listed in \S6 of the report
(\code{fetchMesh}, \code{scheduleRetry}, \code{buildGeometry}, \code{killStub},
the \code{MSG_MESH_DATA} branch, the triplanar helpers, \code{CacheKey} and
\code{CacheKeyHash}, \code{meshCache_}, \code{cacheMutex_},
\code{lastRetryMs_}, the six fetch members, the fetch wakeups, and the
cache-prune loop) while keeping every message constant, the server handlers,
and \code{sendMutex_}; \code{chunkproto::buildRequestMesh} and
\code{buildRequestAllMeshes}; \code{WsClient::sendBinary(const uint8_t*,
size_t)}; \code{WebServer::textureRoot}; \code{SvgShape::rgbToHsvDegrees} and
\code{contains}; \code{SvgPolygonDistanceFunction::localProfilePoint};
\code{SdfScene::createFireDemo}. Delete
\code{vulkan/resources/EditableTexture.{cpp,hpp}} and the include in
\code{VegetationRenderer.hpp}.

Group 6 --- application and widgets. Delete \code{MyApp::dispatchSolidEvents}
and \code{dispatchLiquidEvents}; \code{RenderTargetsWidget::cleanup} (keep the
destructor and \code{destroyLinearTargets}); and every remaining inline
accessor in \S5 except the kept \code{IndirectRenderer} family: the seven
\code{AudioAnalyzer} getters, \code{ParticlePhysicsSim::regionNodeCount},
\code{SvgPolygonDistanceFunction::halfExtents},
\code{MainSceneLoader::shapeIndex}, \code{layerRendererName},
\code{textureTypeName}, \code{MusicWidget::isPlaying},
\code{OctreeExplorerWidget::setExportBlenderCallback},
\code{ImGuiHelpers::openPtr}, \code{DescriptorUpdateStats::updateCalls},
\code{SlotAllocator}'s four helpers, \code{SceneQueues::getGeometryQueue},
\code{CommandBufferState::bindComputePipeline} and
\code{bindGraphicsPipelineUnchecked}, \code{DescriptorBuffer::bindingOffset},
\code{DescriptorBufferHelper::bindingInfo},
\code{SolidRenderer::activeGraphicsPipeline},
\code{RayTracingResources::getRTSetLayout}, the two \code{WaterRenderer}
selectors, \code{IndirectRenderer}'s two size getters, and the three
\code{vk_layouts} helpers. Do \emph{not} delete the other
\code{IndirectRenderer} getters or \code{bindSet0ForFrame}; see \S7 of the
report.

Group 7 --- perf-fix residue. Delete \code{ThreadPool::enqueueDetached} and
\code{notifyOne} (their callers moved to
\code{enqueueDetachedNoNotify}/\code{notifyAll} in \code{26f06f3e}), including
the template definition in \code{ThreadPool.tpp} and the declaration in
\code{ThreadPool.hpp}.

After each group: \code{make -j8 debug}; zero warnings and zero errors. Grep
for every removed symbol. Report per group: files changed, lines removed,
build result, and any symbol that turned out to be live (restore it and say
why). Commit each group separately with
\code{MSG="dead-code(<scope>): ..." make commit}. Do not run the app, do not
run valgrind or callgrind, and do not run \code{pdflatex}.
\end{tcolorbox}

\begin{tcolorbox}[promptbox={Prompt --- short form}]
Delete the residual dead code from \code{docs/dead_code_02.tex} at HEAD
\code{9f0d76cb}, one commit per directory group, verifying with
\code{make -j8 debug} and class-qualified greps. The set is: the 74 methods in
Appendix~A; \code{EditableTexture}, \code{Frustum}, and
\code{OctreeTopologyStats}; the \code{RemoteScene} mesh-fetch cluster (keep
the protocol constants, server handlers, and \code{sendMutex_}); the
\code{ThreadPool} pair left dead by the L8 fix; and the conservative inline
subset in \S5. Keep \code{IndirectRenderer::addMeshSlotted}/\code{uploadSlot},
\code{SceneRenderer::bindSet0ForFrame}, the \code{IndirectRenderer} buffer
getters, \code{GamepadPublisher::aButtonPressed}, and the Y-axis
\code{BoundingBox} setters. Do not touch the \code{sdf/} working tree. No
behaviour changes, no refactors, no app runs, no PDF builds.
\end{tcolorbox}

\appendix
\section{Complete Dead-Method Table}
\label{app:full}

Methods whose symbol is never a relocation target in any project object and
whose name has no call-shaped textual use outside declaration/definition.
``Lines'' is the brace-matched body span. The four \code{RemoteScene} rows are
part of the feature path in \S\ref{sec:remote}; the two \code{Frustum} rows are
part of the whole-unit deletion in \S\ref{sec:units}.

\begin{center}\footnotesize
\begin{longtable}{p{0.44\textwidth}p{0.42\textwidth}r}
\toprule
Location & Symbol & Lines \\
\midrule
\endfirsthead
\toprule
Location & Symbol & Lines \\
\midrule
\endhead
\code{services/AudioAnalyzer.cpp:41--44} & \code{AudioAnalyzer::setBandEdges} & 4 \\
\code{vulkan/resources/MaterialManager.cpp:95--107} & \code{MaterialManager::destroy} & 13 \\
\code{services/OctreeTopologyCompute.cpp:783--793} & \code{OctreeTopologyCompute::stats()} & 11 \\
\code{services/OctreeTopologyCompute.cpp:137--165} & \code{OctreeTopologyCompute::destroy} & 29 \\
\code{services/TextureMixer.cpp:157--163} & \code{TextureMixer::isLayerGenerationPending} & 7 \\
\code{services/TextureMixer.cpp:165--193} & \code{TextureMixer::waitForLayerGeneration} & 29 \\
\code{services/TextureMixer.cpp:32--32} & \code{TextureMixer::getGlobalInstance()} & 1 \\
\code{services/VideoTextureService.cpp:1411--1417} & \code{VideoTextureService::backendCompiledIn()} & 7 \\
\code{services/VideoTextureService.cpp:1437--1453} & \code{VideoTextureService::findAssetDirectory()} & 17 \\
\code{services/VideoTextureService.cpp:1633--1638} & \code{VideoTextureService::isSlotLoaded(slot) const} & 6 \\
\code{services/VideoTextureService.cpp:1640--1640} & \code{VideoTextureService::videoArrayImage() const} & 1 \\
\code{services/VideoTextureService.cpp:1645--1645} & \code{VideoTextureService::videoArrayMips() const} & 1 \\
\code{services/VideoTextureService.cpp:1647--1647} & \code{VideoTextureService::videoArrayHeight() const} & 1 \\
\code{services/VideoTextureService.cpp:1742--1746} & \code{VideoTextureService::slotOf(index) const} & 5 \\
\code{services/VideoTextureService.cpp:1796--1806} & \code{VideoTextureService::togglePlaying(index)} & 11 \\
\code{vulkan/core/VulkanApp.cpp:1348--1373} & \code{VulkanApp::generateMipmaps} & 26 \\
\code{vulkan/core/VulkanApp.cpp:3111--3127} & \code{VulkanApp::recordTrackedLayoutForCommandBuffer} & 17 \\
\code{vulkan/core/VulkanApp.cpp:3503--3539} & \code{VulkanApp::updateDescriptorSet(vector)} & 37 \\
\code{vulkan/core/VulkanResourceManager.cpp:151--151} & \code{VulkanResourceManager::removeFramebuffer} & 1 \\
\code{vulkan/core/VulkanResourceManager.cpp:206--206} & \code{VulkanResourceManager::getSemaphoreMap()} & 1 \\
\code{vulkan/resources/TextureArrayManager.cpp:178--198} & \code{TextureArrayManager::destroy} & 21 \\
\code{vulkan/resources/TextureArrayManager.cpp:465--556} & \code{TextureArrayManager::load} & 92 \\
\code{vulkan/renderer/SceneRenderer.cpp:332--343} & \code{SceneRenderer::removeAllMeshesSolidOnly()} & 12 \\
\code{vulkan/renderer/SceneRenderer.cpp:345--355} & \code{SceneRenderer::removeAllMeshesWaterOnly()} & 11 \\
\code{vulkan/renderer/SceneRenderer.cpp:202--205} & \code{SceneRenderer::layerStateCount()} & 4 \\
\code{vulkan/renderer/sdf/SdfRenderer.cpp:480--486} & \code{SdfRenderer::setSmokeWideContainer} & 7 \\
\code{vulkan/renderer/sdf/SdfRenderer.cpp:322--326} & \code{SdfRenderer::setScene} & 5 \\
\code{vulkan/renderer/sdf/SdfRenderer.cpp:636--638} & \code{SdfRenderer::updateScene()} & 3 \\
\code{vulkan/renderer/sdf/SdfRenderer.cpp:679--684} & \code{SdfRenderer::setMarchParams} & 6 \\
\code{vulkan/renderer/sdf/SdfRenderer.cpp:686--692} & \code{SdfRenderer::setMarchRange} & 7 \\
\code{vulkan/renderer/sdf/SdfRenderer.cpp:714--722} & \code{SdfRenderer::setMarchTiers} & 9 \\
\code{vulkan/renderer/sdf/SdfRenderer.cpp:2142--2146} & \code{SdfRenderer::onSwapchainResized} & 5 \\
\code{vulkan/renderer/video/VideoRenderer.cpp:406--409} & \code{VideoRenderer::clear} & 4 \\
\code{vulkan/renderer/water/WaterBackFaceRenderer.cpp:14--16} & \code{WaterBackFaceRenderer::init} & 3 \\
\code{vulkan/renderer/vegetation/VegetationRenderer.cpp:1003--1020} & \code{VegetationRenderer::setBillboardArrayTextures} & 18 \\
\code{vulkan/renderer/vegetation/ImpostorCapture.cpp:454--508} & \code{ImpostorCapture::capture} & 55 \\
\code{vulkan/renderer/vegetation/ImpostorCapture.cpp:510--519} & \code{ImpostorCapture::captureAll} & 10 \\
\code{vulkan/renderer/vegetation/ImpostorCapture.cpp:521--558} & \code{ImpostorCapture::captureAsync} & 38 \\
\code{vulkan/renderer/vegetation/ImpostorCapture.cpp:1078--1089} & \code{ImpostorCapture::updateTexDescSet} & 12 \\
\code{events/NunchukPublisher.cpp:792--798} & \code{NunchukPublisher::aButtonPressed()} & 7 \\
\code{events/EventManager.cpp:21--25} & \code{EventManager::unsubscribe} & 5 \\
\code{space/Octree.cpp:218--227} & \code{Octree::fetch} & 10 \\
\code{space/Octree.cpp:145--167} & \code{Octree::getNodeAt(pos, level, simplification)} & 23 \\
\code{space/Octree.cpp:169--189} & \code{Octree::getNodeAt(pos, simplification)} & 21 \\
\code{space/OctreeFile.cpp:61--93} & \code{OctreeFile::load} & 33 \\
\code{space/OctreeFile.cpp:95--129} & \code{OctreeFile::save} & 35 \\
\code{utils/GltfWriter.cpp:93--98} & \code{GltfWriter::vertexCount()} & 6 \\
\code{utils/GltfWriter.cpp:100--105} & \code{GltfWriter::triangleCount()} & 6 \\
\code{utils/LocalScene.cpp:297--297} & \code{LocalScene::getOpaqueOctree()} & 1 \\
\code{utils/LocalScene.cpp:298--300} & \code{LocalScene::getOpaqueOctree() const} & 3 \\
\code{utils/LocalScene.cpp:301--305} & \code{LocalScene::getTransparentOctree()} & 5 \\
\code{utils/LocalScene.cpp:306--310} & \code{LocalScene::getTransparentOctree() const} & 5 \\
\code{utils/RemoteScene.cpp:578--600} & \code{RemoteScene::fetchMesh} & 23 \\
\code{utils/RemoteScene.cpp:602--626} & \code{RemoteScene::scheduleRetry} & 25 \\
\code{utils/RemoteScene.cpp:650--688} & \code{RemoteScene::buildGeometry} & 39 \\
\code{utils/RemoteScene.cpp:275--283} & \code{RemoteScene::killStub} & 9 \\
\code{server/ChunkProtocol.cpp:99--101} & \code{buildRequestAllMeshes()} & 3 \\
\code{server/ChunkProtocol.cpp:90--97} & \code{buildRequestMesh} & 8 \\
\code{sdf/SvgShape.cpp:492--505} & \code{SvgShape::rgbToHsvDegrees} & 14 \\
\code{sdf/SvgShape.cpp:507--509} & \code{SvgShape::contains} & 3 \\
\code{sdf/SvgPolygonDistanceFunction.cpp:88--93} & \code{SvgPolygonDistanceFunction::localProfilePoint} & 6 \\
\code{sdf/types/SdfScene.cpp:363--420} & \code{SdfScene::createFireDemo} & 58 \\
\code{math/geometry/BoundingBox.cpp:48--50} & \code{BoundingBox::setMaxX} & 3 \\
\code{math/geometry/BoundingBox.cpp:56--58} & \code{BoundingBox::setMaxZ} & 3 \\
\code{math/geometry/BoundingBox.cpp:60--62} & \code{BoundingBox::setMinX} & 3 \\
\code{math/geometry/BoundingBox.cpp:68--70} & \code{BoundingBox::setMinZ} & 3 \\
\code{math/geometry/BoundingCube.cpp:59--62} & \code{BoundingCube::getChildCenter} & 4 \\
\code{math/geometry/BoundingCube.cpp:65--67} & \code{BoundingCube::intersects(Ray)} & 3 \\
\code{math/geometry/BoundingCube.cpp:69--71} & \code{BoundingCube::setMax} & 3 \\
\code{math/geometry/Frustum.cpp:4--14} & \code{Frustum::Frustum} & 11 \\
\code{math/geometry/Frustum.cpp:16--50} & \code{Frustum::test} & 35 \\
\code{MyApp.cpp:5195--5197} & \code{MyApp::dispatchSolidEvents()} & 3 \\
\code{MyApp.cpp:5199--5206} & \code{MyApp::dispatchLiquidEvents()} & 8 \\
\code{widgets/RenderTargetsWidget.cpp:662--781} & \code{RenderTargetsWidget::cleanup()} & 120 \\
\midrule
\textbf{Total} & \textbf{74 methods} & \textbf{1{,}094} \\
\bottomrule
\end{longtable}
\end{center}

\end{document}
